\documentclass[12pt,a4paper]{report}

\usepackage[utf8]{inputenc}
\usepackage[T1]{fontenc}
\usepackage[french]{babel}
\usepackage{graphicx}
\usepackage{amsmath,amssymb}
\usepackage{booktabs}
\usepackage{geometry}
\usepackage{hyperref}
\usepackage{array}
\usepackage{longtable}
\usepackage{float}
\usepackage{xcolor}
\usepackage{setspace}
\usepackage{tcolorbox}

\geometry{margin=2.5cm}
\setlength{\parindent}{0.6cm}
\setlength{\parskip}{0.25em}
\onehalfspacing

% Couleurs pour le cadre de la page de garde
\definecolor{lightGray}{RGB}{240,240,240}
\definecolor{primaryBlue}{RGB}{0,102,204}

\newcommand{\captureplaceholder}[2]{
	\begin{figure}[htbp]
		\centering
		\fbox{\begin{minipage}[c][5cm][c]{0.88\textwidth}
			\centering
			\textbf{Capture d'écran à insérer}\\[0.4cm]
			#1\\[0.2cm]
			\small #2
		\end{minipage}}
		\caption{#1}
	\end{figure}
}

\newcommand{\screenshotpage}[2]{
\clearpage
\begin{figure}[p]
\centering
\fbox{\begin{minipage}[c][14cm][c]{0.9\textwidth}
\centering
\Large \textbf{Emplacement réservé à une capture d'écran}\\[0.8cm]
\normalsize #1\\[0.5cm]
\small #2
\end{minipage}}
\caption{#1}
\end{figure}
\clearpage
}

\begin{document}

\begin{titlepage}
	\centering

	\vspace*{1cm}
	% Remplacer les chemins ci-dessous par les logos disponibles.
	\begin{minipage}{0.45\textwidth}
		\centering
		\includegraphics[height=2cm]{emsi.png}
	\end{minipage}
	\hfill
	\begin{minipage}{0.45\textwidth}
		\centering
		\includegraphics[height=2cm]{logo (2).png}
	\end{minipage}

	\vfill

	{\Huge \textbf{Rapport de Projet de Fin de Stage}\par}
	\vspace{0.8cm}

	\begin{tcolorbox}[colback=lightGray,colframe=primaryBlue,arc=2mm,boxrule=1pt,width=0.94\textwidth]
		\centering
		{\Huge \textbf{Démonstrateur de prévision horaire de la demande électrique}\par}
	\end{tcolorbox}

	\vspace{2.0cm}
	{\large Entreprise d'accueil : \textbf{Mogador Technologie Industrielle}\par}
	\vspace{0.3cm}
	{\large Période de stage : \textbf{Du 10 août au 25 septembre 2026}\par}

	\vfill
	\begin{center}
		\large
		\textbf{Étudiant :} Elhannak Badr\\[0.3cm]
		\textbf{Encadrant pédagogique :} Saddouk Lamyae\\[0.3cm]
		\textbf{Encadrant entreprise :} Khatar Zakaria
	\end{center}

	\vfill

	{\large Année universitaire : 2025--2026\par}
\end{titlepage}


\chapter*{Remerciements}
\addcontentsline{toc}{chapter}{Remerciements}
Je tiens tout d’abord à exprimer ma profonde gratitude à l’ensemble des personnes qui ont contribué, de près ou de loin, à la réalisation de ce projet de fin de stage. Cette expérience professionnelle a constitué une étape importante dans mon parcours académique et m’a permis de mettre en pratique les connaissances acquises au cours de ma formation, tout en développant de nouvelles compétences dans un environnement professionnel.

J’adresse tout particulièrement mes sincères remerciements à l’ensemble de l’équipe de **Mogador Technologie Industrielle** pour son accueil chaleureux, sa disponibilité et son accompagnement tout au long de mon stage. Je remercie également les collaborateurs qui ont pris le temps de répondre à mes questions, de partager leurs connaissances et de me fournir les informations nécessaires à la compréhension du fonctionnement de l’entreprise, de son environnement professionnel ainsi que des différentes problématiques liées au projet.

Je souhaite également remercier mon **encadrant en entreprise** pour son accompagnement, ses conseils et ses orientations tout au long de cette expérience. Son suivi et ses remarques m’ont permis de mieux comprendre les exigences du milieu professionnel, d’améliorer ma méthode de travail et d’avancer progressivement dans la réalisation des différentes étapes du projet.

Mes remerciements s’adressent également à mon **encadrant pédagogique**, pour son suivi, ses conseils et ses recommandations, qui m’ont permis de structurer mon travail et de maintenir une démarche cohérente avec les objectifs pédagogiques de cette période de stage.



Enfin, je remercie toutes les personnes, enseignants, encadrants, collaborateurs et collègues, qui ont contribué, directement ou indirectement, au bon déroulement de ce stage et à l’enrichissement de mon expérience. Leur disponibilité, leurs conseils et leur soutien ont largement participé à la réussite de cette expérience professionnelle et à mon développement personnel et professionnel.


\chapter*{Résumé}
\addcontentsline{toc}{chapter}{Résumé}

Ce projet consiste à concevoir un démonstrateur local de prévision horaire de
la demande électrique. La solution exploite un historique de données horaires
comprenant la demande, la température et l'humidité. Les données utilisées
sont synthétiques et couvrent la période allant de janvier 2020 à décembre
2024. Après une étape d'ingestion et de contrôle qualité, des variables
calendaires, météorologiques, de retard et des statistiques glissantes sont
construites afin d'alimenter un modèle de régression XGBoost. L'évaluation
respecte l'ordre chronologique : les données jusqu'à 2022 servent à
l'entraînement, l'année 2023 à la sélection et l'année 2024 au test final.
Le modèle retenu obtient sur le test 2024 un MAE de 103,72 MW, un RMSE de
130,22 MW, un coefficient de détermination $R^2$ de 0,9280 et un WAPE de
3,55\%. Le système est exposé par une API FastAPI et visualisé dans une
interface React proposant l'exploration des données, l'analyse de la demande,
l'évaluation du modèle et la prévision future récursive. Le résultat constitue
un prototype académique et local ; il ne prétend pas être une solution
industrielle validée sur des mesures électriques réelles.

\textbf{Mots-clés :} demande électrique, prévision horaire, XGBoost, FastAPI,
React, séries temporelles.

\chapter*{Abstract}
\addcontentsline{toc}{chapter}{Abstract}

This project aims to build a local demonstrator for hourly electricity demand
forecasting. The solution uses a historical hourly dataset containing demand,
temperature and humidity. The data are synthetic and cover the period from
January 2020 to December 2024. After ingestion and quality checks, calendar,
weather, lagged-demand and rolling-statistics features are generated for an
XGBoost regression model. The evaluation follows the chronological order of
the observations: data up to 2022 are used for training, 2023 is used for
model selection and early stopping, and 2024 is reserved for the final test.
On the 2024 test set, the selected model achieves a MAE of 103.72 MW, an RMSE
of 130.22 MW, an $R^2$ score of 0.9280 and a WAPE of 3.55\%. The system is
served through a FastAPI backend and displayed through a React interface that
provides data exploration, demand analysis, model evaluation and recursive
future forecasting. The result is an academic local prototype. It has not
been validated as an operational industrial solution on real electricity
measurements.

\textbf{Keywords:} electricity demand, hourly forecasting, XGBoost, FastAPI,
React, time series.

% Ajouter éventuellement une version arabe sur le même modèle.

\tableofcontents
\addcontentsline{toc}{chapter}{Table des matières}

\listoffigures
\addcontentsline{toc}{chapter}{Liste des figures}

\listoftables
\addcontentsline{toc}{chapter}{Liste des tableaux}

\chapter*{Liste des abréviations}
\addcontentsline{toc}{chapter}{Liste des abréviations}

\begin{longtable}{p{0.22\textwidth} p{0.68\textwidth}}
\caption{Liste des abréviations}\label{tab:abreviations}\\
\toprule
Abréviation & Signification \\
\midrule
\endfirsthead
\toprule
Abréviation & Signification \\
\midrule
\endhead
API & Application Programming Interface \\
BI & Business Intelligence \\
CSV & Comma-Separated Values \\
ETL & Extract, Transform and Load \\
F1 & F-measure, moyenne harmonique de la précision et du rappel \\
IA & Intelligence artificielle \\
JSON & JavaScript Object Notation \\
MAE & Mean Absolute Error, erreur absolue moyenne \\
ML & Machine Learning, apprentissage automatique \\
MW & Megawatt, unité de puissance \\
RMSE & Root Mean Squared Error, racine de l'erreur quadratique moyenne \\
UML & Unified Modeling Language \\
URL & Uniform Resource Locator \\
WAPE & Weighted Absolute Percentage Error \\
XLSX & Format de classeur Microsoft Excel \\
\bottomrule
\end{longtable}

\chapter{Introduction générale}

\section{Contexte général du stage}

Ce stage s'inscrit dans le cadre de la formation d'ingénieur et porte sur la
mise en place d'un outil d'analyse et de prévision de la demande électrique.
Il permet de mobiliser des compétences en traitement des données,
apprentissage automatique, développement backend et visualisation web.

Le projet est réalisé dans le contexte de Mogador Technologie Industrielle,
entreprise marocaine située au N 55, boulevard Mohamed Bouziane, Sidi
Othmane, Casablanca. L'entreprise opère dans le secteur des installations
électriques et hydrauliques ainsi que dans la coordination d'entreprises.
Dans cet environnement, une meilleure lecture de la demande électrique peut
faciliter l'analyse des historiques et la préparation d'outils d'aide à la
décision.

\section{Problématique posée}

La problématique est la suivante : comment exploiter un historique de demande
électrique et de variables météorologiques pour produire une prévision horaire
lisible, évaluable et accessible à travers une interface web ?

La difficulté provient du caractère temporel des données, de la nécessité de
respecter la causalité lors de la création des variables et de la différence
entre une prévision historique évaluée avec des observations disponibles et
une prévision future entièrement récursive. Le projet doit donc fournir une
chaîne cohérente tout en distinguant clairement le prototype technique d'une
solution opérationnelle.

\section{Objectifs du projet}

L'objectif principal est de développer un démonstrateur local capable
d'explorer la demande électrique, d'évaluer un modèle de prévision et de
générer une prévision future sur un horizon de une à 168 heures. Les objectifs
spécifiques sont les suivants :

\begin{itemize}
	\item contrôler et préparer les données horaires ;
	\item construire des variables calendaires, météorologiques et historiques ;
	\item entraîner et comparer un modèle XGBoost à des baselines naïves ;
	\item exposer les traitements par une API FastAPI ;
	\item développer une interface React pour l'analyse et la prévision ;
	\item évaluer les performances et documenter les limites du prototype.
\end{itemize}

\section{Plan du rapport}

Le premier chapitre présente Mogador Technologie Industrielle, le contexte
du projet et le cahier des charges. Le deuxième chapitre expose l'état de
l'art et justifie le choix des technologies. Le troisième chapitre décrit la
conception, l'architecture et la modélisation de la solution. Le quatrième
chapitre présente la réalisation, l'environnement technique et les tests. Le
cinquième chapitre analyse les résultats et leurs limites. Enfin, la
conclusion générale dresse le bilan du travail et présente les perspectives.

\chapter{Contexte et Cahier des Charges}

\section{Présentation de l'entreprise d'accueil}

Mogador Technologie Industrielle est une entreprise basée au Maroc. Son siège
est situé au N 55, boulevard Mohamed Bouziane, Sidi Othmane, Casablanca.
L'entreprise opère dans le secteur des installations électriques,
hydrauliques et de la coordination d'entreprises.

Son activité la place au croisement de plusieurs intervenants techniques et
organisationnels. Les installations électriques et hydrauliques nécessitent en
particulier une bonne coordination des opérations, un suivi des équipements
et une capacité à exploiter les informations disponibles pour préparer les
décisions techniques.

\begin{figure}[H]
    \centering
    \includegraphics[width=0.6\textwidth]{M1.jpeg}
    \caption{Présentation de Mogador Technologie Industrielle}
    \label{fig:}
\end{figure}

\section{Contexte métier du projet}

Le projet concerne l'exploitation de données de demande électrique dans un
contexte d'installations et de coordination techniques. Les acteurs concernés
sont les responsables qui consultent les indicateurs, les utilisateurs qui
analysent l'évolution de la demande et les profils techniques qui peuvent
utiliser les prévisions comme support d'analyse.

Le démonstrateur ne commande pas directement une installation. Il fournit une
lecture structurée de l'historique et une estimation de la demande future afin
de faciliter l'analyse et l'expérimentation.

\section{Analyse de l'existant}

Avant la réalisation du démonstrateur, l'analyse pouvait être effectuée à
partir de fichiers de données et de traitements séparés. Cette organisation
rend plus difficile la vérification de la qualité des données, la comparaison
des modèles et la consultation homogène des résultats.

La solution développée regroupe l'ingestion, la création des variables,
l'évaluation du modèle, l'exposition des résultats et leur visualisation dans
une même chaîne logicielle.

\subsection{Problèmes et limites de la solution actuelle}

Les principales limites identifiées sont l'absence d'une interface unifiée,
la difficulté à reproduire les traitements manuels et le manque de visibilité
sur les métriques de prévision. Le projet répond à ces limites par une API
structurée, des contrôles de qualité et un tableau de bord composé de pages
d'analyse dédiées.

\section{Expression des besoins}

\subsection{Besoins fonctionnels}

Les besoins fonctionnels retenus sont les suivants :

\begin{itemize}
	\item charger un fichier de données CSV ou XLSX ;
	\item vérifier les doublons, les trous horaires et les valeurs manquantes ;
	\item présenter les statistiques générales et l'évolution de la demande ;
	\item analyser les effets de l'heure, du mois, du week-end et de la météo ;
	\item afficher les métriques du modèle et les comparer aux baselines ;
	\item consulter une prédiction historique à une date donnée ;
	\item générer une prévision future de 1 à 168 heures ;
	\item afficher les informations dans une interface React.
\end{itemize}

\subsection{Besoins non fonctionnels}

Les besoins non fonctionnels sont la reproductibilité des traitements,
l'absence de fuite temporelle, la lisibilité des graphiques, la validation des
paramètres reçus par l'API et la séparation claire entre backend et frontend.
La solution doit également rester simple à exécuter localement et fournir des
messages d'erreur compréhensibles lorsque les données ou les paramètres sont
invalides.

\section{Contraintes techniques et organisationnelles}

Le projet est contraint par l'utilisation de données synthétiques, l'absence
d'un fournisseur météo externe et le caractère local du démonstrateur. Le
	fichier de données doit respecter le contrat de colonnes \texttt{Timestamp},
\texttt{Demand}, \texttt{Temperature} et \texttt{Humidity}. La demande cible n'est pas imputée,
alors que les variables météorologiques peuvent être propagées vers l'avant.

Une autre contrainte est la prévision récursive : au-delà de la première
heure, les prédictions précédentes alimentent les étapes suivantes. Les
résultats multi-horizon doivent donc être considérés comme expérimentaux et
non comme une garantie opérationnelle.

\section{Périmètre du projet}

Le périmètre fonctionnel comprend l'analyse d'un historique horaire, la
préparation des données, la construction de variables explicatives,
l'entraînement d'un modèle de régression, l'exposition des résultats par une
API et leur consultation dans une interface web. Le périmètre comprend aussi
la comparaison avec des méthodes de référence et la génération d'une
prévision récursive limitée à 168 heures dans l'application.

En revanche, le projet ne couvre pas le pilotage automatique d'une
installation, la facturation, la détection de défauts électriques, la
prévision météorologique elle-même ou le déploiement sur une infrastructure
de production. Cette délimitation est importante : le prototype démontre la
faisabilité technique d'une chaîne de prévision, mais ne remplace pas une
étude industrielle complète.

\section{Critères d'acceptation}

Les critères d'acceptation retenus sont présentés dans le tableau
\ref{tab:criteres-acceptation}. Ils permettent de vérifier que chaque besoin
possède un résultat observable et un moyen de validation.

\begin{longtable}{p{0.27\textwidth} p{0.38\textwidth} p{0.22\textwidth}}
\caption{Critères d'acceptation du projet}\label{tab:criteres-acceptation}\\
\toprule
Critère & Résultat attendu & Vérification \\
\midrule
\endfirsthead
\toprule
Critère & Résultat attendu & Vérification \\
\midrule
\endhead
Qualité des données & Les doublons, trous et valeurs manquantes sont identifiés. & Tests d'ingestion \\
Variables causales & Les features de demande utilisent uniquement le passé. & Tests de features \\
Modèle & Le modèle produit une prédiction numérique. & Test de prédiction \\
Évaluation & Les métriques sont calculées sur un test chronologique. & Tests de modélisation \\
API & Les endpoints répondent avec des données structurées. & Tests FastAPI \\
Interface & Les quatre pages principales s'affichent. & Tests frontend \\
Prévision future & Un horizon valide produit une série horaire. & Test récursif \\
\bottomrule
\end{longtable}

\section{Risques du projet}

Plusieurs risques ont été identifiés dès la conception. Le premier est le
risque de fuite temporelle : une variable calculée avec l'observation cible
ou une observation future rendrait les performances artificiellement
optimistes. Il est réduit par le décalage des historiques et par la séparation
chronologique des jeux de données.

Le deuxième risque est la dégradation des prévisions récursives. Lorsque les
prédictions sont réinjectées, une erreur à une étape peut influencer les
étapes suivantes. Le rapport distingue donc les performances historiques des
prévisions réellement futures.

Le troisième risque est lié à la qualité et à l'origine des données. Le jeu
utilisé est synthétique ; les résultats ne doivent pas être généralisés à une
installation réelle sans nouvelle campagne de validation.

\begin{table}[htbp]
\centering
\caption{Matrice synthétique des risques}
\label{tab:risques}
\begin{tabular}{p{0.28\textwidth} p{0.16\textwidth} p{0.42\textwidth}}
\toprule
Risque & Niveau & Mesure de réduction \\
\midrule
Fuite de données temporelles & Élevé & Décalage des features et séparation par année \\
Données synthétiques & Élevé & Validation future sur des mesures réelles \\
Erreur récursive & Élevé & Limiter l'horizon et afficher une mise en garde \\
Météo indisponible & Moyen & Utiliser un proxy historique explicitement documenté \\
Incompatibilité de schéma & Moyen & Vérifier les colonnes et l'ordre des features \\
\bottomrule
\end{tabular}
\end{table}

\chapter{État de l'Art}

\section{Technologies et frameworks étudiés}

La solution s'appuie sur Python pour l'ingestion, la préparation et la
modélisation des données. Pandas et NumPy sont utilisés pour les traitements
tabulaires et numériques. Scikit-learn intervient dans l'évaluation et le
calcul des métriques, tandis que XGBoost fournit le modèle de régression.
Joblib permet de charger le modèle pré-entraîné et OpenPyXL prend en charge
la lecture des fichiers Excel.

Le backend est construit avec FastAPI et exécuté avec Uvicorn. Le frontend
est une application React 19 construite avec Vite. Les graphiques sont
réalisés avec Recharts et les tests frontend utilisent Vitest et Testing
Library.

\section{Choix technologique justifié}

XGBoost a été retenu car il convient aux données tabulaires, permet de
modéliser des relations non linéaires et offre un compromis intéressant entre
performance et temps de calcul. FastAPI fournit une API légère et typée pour
relier le modèle à l'interface utilisateur. React et Recharts permettent de
présenter plusieurs vues analytiques sans mélanger la logique métier et la
présentation.

Ces choix restent adaptés à un prototype local. Une mise en production
nécessiterait notamment une validation sur des données réelles, une gestion
des utilisateurs, une supervision et une stratégie de déploiement.

\begin{table}[htbp]
	\centering
	\caption{Comparaison des alternatives technologiques}
	\label{tab:comparaison-technologies}
	\begin{tabular}{p{0.25\textwidth} p{0.3\textwidth} p{0.3\textwidth}}
		\toprule
		Alternative & Avantages & Inconvénients \\
		\midrule
		Traitement manuel & Simple à démarrer & Peu reproductible \\
		Modèle naïf & Référence facile à interpréter & Précision limitée \\
		XGBoost + API web & Bon pour les données tabulaires et consultable & Nécessite une chaîne de données fiable \\
		\bottomrule
	\end{tabular}
\end{table}

% Pour un projet d'IA : présenter les modèles et algorithmes candidats.
% Pour un projet de BI : présenter les outils ETL, entrepôts et outils de visualisation.

\section{Prévision de séries temporelles}

Une série temporelle est une suite d'observations indexées par le temps. Dans
le cas de la demande électrique, chaque observation correspond à une heure et
la variable cible est la puissance demandée. La prévision consiste à estimer
une valeur future à partir des observations disponibles avant cette heure.

La structure de la demande combine généralement plusieurs composantes : une
tendance de long terme, une saisonnalité journalière, une saisonnalité
hebdomadaire, des effets calendaires et une composante irrégulière. Le projet
représente ces composantes à l'aide de variables temporelles et de retards de
la demande plutôt qu'avec un modèle statistique classique unique.

Pour éviter une discontinuité artificielle entre 23 heures et 0 heure, l'heure
est également représentée par des fonctions trigonométriques :

\begin{equation}
hour\_sin_t = \sin\left(2\pi \frac{hour_t}{24}\right),
\qquad
hour\_cos_t = \cos\left(2\pi \frac{hour_t}{24}\right).
\end{equation}

La même logique est utilisée pour le jour de la semaine. Les variables
cycliques permettent au modèle de comprendre que 23 heures et 0 heure sont
voisines, contrairement à une représentation numérique brute.

\section{Méthodes de référence}

Une baseline est indispensable pour savoir si un modèle d'apprentissage
apporte réellement une amélioration. La baseline à 24 heures suppose que la
demande observée la veille à la même heure constitue une bonne estimation. La
baseline à 168 heures reprend la demande de la semaine précédente.

Ces méthodes sont simples, transparentes et difficiles à battre lorsque la
saisonnalité est très régulière. Elles présentent toutefois des limites en
cas de variation météorologique, de changement de comportement ou de journée
exceptionnelle. Le modèle XGBoost cherche à combiner ces répétitions avec les
variables calendaires et météorologiques.

\section{Apprentissage par gradient boosting}

XGBoost construit une somme séquentielle d'arbres de décision. À chaque
itération, un nouvel arbre tente de corriger les erreurs résiduelles des
arbres précédents. Pour une régression, la prédiction peut être représentée
schématiquement par :

\begin{equation}
\hat{y}_t = \sum_{k=1}^{K} f_k(x_t),
\end{equation}

où $K$ est le nombre d'arbres et $f_k$ un arbre appartenant à l'espace des
modèles faibles. Le taux d'apprentissage contrôle la contribution de chaque
arbre. La profondeur, le sous-échantillonnage et les régularisations limitent
le surapprentissage.

Le choix de XGBoost est cohérent avec un jeu de données tabulaire enrichi de
variables temporelles. Il ne nécessite pas de normaliser toutes les variables
et peut représenter des relations non linéaires entre température, calendrier
et demande.

\section{Comparaison avec des alternatives}

Une moyenne mobile est facile à calculer mais réagit lentement aux variations
rapides. Une régression linéaire est interprétable mais peut être trop
restrictive pour les interactions entre variables. Les réseaux récurrents ou
les transformeurs pourraient représenter des dépendances complexes, mais ils
nécessitent davantage de données, de réglages et de ressources.

Dans ce projet académique, XGBoost offre un compromis raisonnable entre
complexité, vitesse d'entraînement, interprétabilité par importance des
variables et qualité des résultats.

\begin{table}[htbp]
\centering
\caption{Comparaison élargie des approches de prévision}
\label{tab:approches-prevision}
\begin{tabular}{p{0.22\textwidth} p{0.26\textwidth} p{0.22\textwidth} p{0.2\textwidth}}
\toprule
Approche & Principe & Atout & Limite \\
\midrule
Naïve 24 h & Répétition journalière & Très simple & Ignore les changements \\
Naïve 168 h & Répétition hebdomadaire & Bonne référence saisonnière & Réagit peu aux événements \\
Régression linéaire & Relation additive & Interprétable & Relations limitées \\
XGBoost & Ensemble d'arbres & Non linéaire et efficace & Dépend des features \\
Réseau profond & Représentation séquentielle & Flexible & Plus coûteux à régler \\
\bottomrule
\end{tabular}
\end{table}

\chapter{Conception et Architecture}

\section{Méthodologie de gestion de projet}

Le projet a été mené selon une démarche itérative. Les travaux ont été
organisés autour de plusieurs étapes : compréhension du besoin, exploration
des données, construction des variables, entraînement et comparaison des
modèles, développement de l'API, développement de l'interface et validation.

Cette organisation est proche d'une démarche Agile, car chaque étape produit
un résultat vérifiable et peut conduire à des ajustements. Elle est adaptée à
un prototype de recherche appliquée dont les besoins techniques sont précisés
progressivement.

\section{Architecture globale de la solution}

L'architecture globale suit le flux suivant : le fichier CSV ou XLSX est chargé
par le module d'ingestion, contrôlé puis transformé en variables utilisables
par le modèle XGBoost. Le modèle pré-entraîné est chargé avec Joblib. FastAPI
expose ensuite les données, les métriques et les prédictions sous forme
d'endpoints HTTP. Enfin, l'application React consomme ces endpoints et
présente les résultats avec des graphiques Recharts.

L'architecture est organisée en couches. La couche données lit le fichier et
effectue les contrôles de structure. La couche features enrichit les lignes
avec les variables causales. La couche modèle charge le package XGBoost et
calcule les métriques. La couche API transforme ces résultats en réponses
JSON. Enfin, la couche frontend sélectionne les routes et rend les graphiques.

Cette séparation limite le couplage. Une modification de l'affichage ne doit
pas modifier les règles de calcul des features. De même, une évolution du
modèle doit conserver le contrat de réponse attendu par les pages React.

\begin{table}[htbp]
\centering
\caption{Responsabilités des couches de la solution}
\label{tab:couches-architecture}
\begin{tabular}{p{0.22\textwidth} p{0.35\textwidth} p{0.28\textwidth}}
\toprule
Couche & Responsabilité & Modules principaux \\
\midrule
Données & Lecture et contrôle des fichiers & ingestion.py \\
Features & Création des variables causales & features.py \\
Modèle & Entraînement, métriques et artefact & modeling.py \\
API & Validation et exposition HTTP & main.py \\
Interface & Navigation et visualisation & App.jsx, pages React \\
\bottomrule
\end{tabular}
\end{table}

\begin{figure}[H]
    \centering
    \includegraphics[width=1.0\textwidth]{M2.jpeg}
    \caption{Architecture globale de la solution}
    \label{fig:}
\end{figure}

\section{Modélisation et Conception}

La conception repose sur une séparation entre les données, la logique de
modélisation et la présentation. Le backend regroupe l'ingestion, les
features, l'évaluation et les routes API. Le frontend est organisé en pages
correspondant aux besoins d'analyse : vue d'ensemble, analyse de la demande,
performance du modèle et prévision.

Le projet ne contient pas de diagrammes UML formels dans son implémentation
actuelle. Les éléments de conception ci-dessous décrivent donc les interactions
fonctionnelles utiles au démonstrateur.

\section{Contrat des données}

Le fichier d'entrée doit contenir quatre colonnes principales. La colonne
\texttt{Timestamp} identifie l'instant d'observation. La colonne
\texttt{Demand} contient la cible exprimée en mégawatts. Les colonnes
\texttt{Temperature} et \texttt{Humidity} apportent le contexte
météorologique. Le contrôle du schéma est réalisé avant le calcul des
features afin d'éviter qu'une erreur de nom de colonne ne soit découverte
seulement au moment de l'entraînement.

Le jeu initial contient 43 848 lignes et couvre janvier 2020 à décembre 2024.
Après retrait des lignes entièrement vides, 43 844 timestamps sont utilisables
par l'application. Les observations sont triées par date et la fréquence
horaire est imposée. Cette étape rend explicites les trous éventuels au lieu
de les masquer dans un index irrégulier.

\begin{longtable}{p{0.24\textwidth} p{0.23\textwidth} p{0.4\textwidth}}
\caption{Dictionnaire des données d'entrée}\label{tab:dictionnaire-donnees}\\
\toprule
Colonne & Type attendu & Description \\
\midrule
\endfirsthead
\toprule
Colonne & Type attendu & Description \\
\midrule
\endhead
\texttt{Timestamp} & Date-heure & Index temporel horaire \\
\texttt{Demand} & Numérique & Demande électrique cible en MW \\
\texttt{Temperature} & Numérique & Température associée à l'observation \\
\texttt{Humidity} & Numérique & Humidité associée à l'observation \\
\bottomrule
\end{longtable}

\section{Catalogue des variables construites}

Les variables calendaires décrivent l'instant sans utiliser la cible. Elles
comprennent l'heure, le jour de la semaine, le mois, l'année, le jour de
l'année, le trimestre et la semaine ISO. Les indicateurs de week-end et de
jour férié marocain ajoutent une information qualitative.

Les variables de retard représentent la demande passée. Les retards de 1, 2,
3, 6 et 12 heures capturent la dynamique courte. Les retards de 24 et
48 heures représentent la répétition journalière. Les retards de 72, 168 et
336 heures donnent au modèle une mémoire de plusieurs jours et semaines.

Les statistiques glissantes résument la demande passée. Pour une fenêtre $w$,
la moyenne est calculée par :

\begin{equation}
\overline{D}_{t,w} = \frac{1}{w}\sum_{i=1}^{w}D_{t-i}.
\end{equation}

L'indice $t-i$ garantit que la demande de l'heure à prédire n'est pas utilisée.
Le même principe est appliqué à l'écart-type, qui mesure la volatilité récente.

\begin{longtable}{p{0.3\textwidth} p{0.26\textwidth} p{0.28\textwidth}}
\caption{Catalogue des familles de features}\label{tab:catalogue-features}\\
\toprule
Famille & Exemples & Rôle \\
\midrule
\endfirsthead
\toprule
Famille & Exemples & Rôle \\
\midrule
\endhead
Calendrier & hour, month, quarter & Décrire la position temporelle \\
Cyclique & hour\_sin, hour\_cos & Représenter les cycles \\
Événements & is\_weekend, is\_holiday & Distinguer les jours particuliers \\
Météo & Temperature, Humidity & Représenter le contexte extérieur \\
Retards & lag\_1h, lag\_24h, lag\_168h & Exploiter l'historique de demande \\
Fenêtres & mean\_24h, std\_168h & Résumer niveau et variabilité \\
\bottomrule
\end{longtable}

% Pour un projet de BI : inclure le modèle dimensionnel et le pipeline ETL.
% Pour un projet d'IA : inclure le pipeline ML/DL, la préparation des données et le modèle choisi.
% Pour un projet de développement : inclure les diagrammes UML.

\subsection{Diagrammes de cas d'utilisation}

L'utilisateur consulte les informations générales, analyse les facteurs liés
à la demande, compare les modèles et demande une prévision historique ou
future. L'API joue le rôle d'intermédiaire entre l'utilisateur et les données
préparées ; elle valide notamment les dates, les horizons et les tableaux
météorologiques reçus.

\begin{figure}[H]
    \centering
    \includegraphics[width=0.6\textwidth]{m3.jpeg}
    \caption{Diagrammes de cas d’utilisation}
    \label{fig:}
\end{figure}

\subsection{Diagrammes de séquence}

Pour une prévision future, l'utilisateur choisit un horizon compris entre 1
et 168 heures et peut fournir des valeurs de température et d'humidité. Le
frontend envoie une requête POST à l'endpoint \texttt{/api/forecast/future}. Le
backend construit les variables de la première heure, calcule la prédiction,
puis réutilise cette prédiction dans les étapes suivantes. La réponse est
retournée au frontend pour affichage.

\begin{figure}[H]
    \centering
    \includegraphics[width=0.6\textwidth]{m4.jpeg}
    \caption{Diagramme de séquence de la prévision}
    \label{fig:}
\end{figure}

\subsection{Diagramme de classes}

Les principaux composants sont les modules d'ingestion, de création des
features, de modélisation et d'API côté backend, ainsi que les pages et
composants de visualisation côté frontend. Les données échangées sont des
objets JSON contenant les métadonnées, les séries temporelles, les métriques
et les prédictions.

\begin{figure}[H]
    \centering
    \includegraphics[width=0.6\textwidth]{m5.jpeg}
    \caption{Diagramme de classes}
    \label{fig:}
\end{figure}

\subsection{Pipeline d'apprentissage automatique}

Le pipeline ML suit les étapes suivantes : chargement du fichier, contrôle
qualité, préparation temporelle, création des variables causales, séparation
chronologique, entraînement, sélection, évaluation et sauvegarde du modèle.
Le schéma doit montrer que les données de test 2024 restent séparées de
l'entraînement et de l'early stopping réalisé sur 2023.

\begin{figure}[H]

    \centering

    \includegraphics[width=0.9\textwidth]{m6.png}

    \caption{Pipline ML complete}

    \label{fig:}

\end{figure}

\subsection{Conception de l'interface utilisateur}

La conception UI est organisée autour de quatre parcours : vue d'ensemble,
analyse de la demande, performance du modèle et prévision. Chaque écran doit
présenter une hiérarchie visuelle claire, des indicateurs lisibles et des
graphiques permettant de comparer les valeurs. Les actions de l'utilisateur
doivent rester cohérentes avec les validations de l'API, notamment pour les
dates et l'horizon futur.

\begin{figure}[H]

    \centering

    \includegraphics[width=1.1\textwidth]{Maquette_LoFi_Dashboard.png}

    \caption{ Maquette et conception UI}

    \label{fig:}

\end{figure}

\begin{figure}[H]
    \centering
    \includegraphics[width=1.0\textwidth]{dashboard_gen.jpeg}
    \caption{nterface React finale}
    \label{fig:}
\end{figure}

\section{Conception de la prévision récursive}

La prévision historique et la prévision future ne disposent pas des mêmes
informations. Dans l’évaluation 2024, les retards peuvent être calculés à
partir de l’historique observé. En prévision future, seule l’histoire jusqu’à
la dernière observation est connue. La fonction récursive utilise donc la
prédiction produite à l’instant $t$ comme une observation pour calculer
certaines features de $t+1$.

Pour les variables météorologiques futures, l’application peut recevoir des
tableaux de température et d’humidité. En leur absence, elle utilise un proxy
basé sur les moyennes historiques par mois et par heure. Cette solution
permet de démontrer le mécanisme, mais ne constitue pas une prévision
météorologique réelle.

Le contrôle de l’horizon accepte une valeur comprise entre 1 et 168 heures.
Au-delà de cette limite, l’incertitude augmente et les résultats doivent être
interprétés avec prudence. La date de départ doit également correspondre à
l’heure suivant la dernière observation.

\chapter{Réalisation et Implémentation}

\section{Environnement de développement}

L'environnement de développement utilise Python pour le backend et JavaScript
pour le frontend. Les principaux outils et bibliothèques sont :

\begin{itemize}
	\item FastAPI et Uvicorn pour l'API ;
	\item Pandas, NumPy et OpenPyXL pour les données ;
	\item XGBoost, scikit-learn et Joblib pour la modélisation ;
	\item la bibliothèque \texttt{holidays} pour les jours fériés marocains ;
	\item React 19, Vite et Recharts pour l'interface ;
	\item Pytest, Vitest et Testing Library pour les tests.
\end{itemize}

Les versions exactes sont définies dans \texttt{requirements.txt} et
\texttt{frontend/package.json}.

\section{Organisation du code source}

Le backend est séparé en modules spécialisés. Le module d'ingestion charge le
fichier, normalise les timestamps et produit les contrôles de qualité. Le
module de features calcule les variables temporelles, météo et historiques.
Le module de modélisation regroupe la sélection, l'évaluation et la
comparaison des configurations. Le module principal initialise l'application
FastAPI et expose les routes.

Le frontend suit une organisation par pages. La page Overview présente les
indicateurs de volume et de période. Demand Analysis rassemble les analyses
par heure, mois, type de jour et température. Model Performance présente les
métriques et les erreurs. Forecast permet de consulter une prédiction
historique ou de demander un horizon futur.

Cette organisation facilite la maintenance, car chaque page possède une
responsabilité claire. Les composants réutilisables, tels que les métriques et
les tableaux, évitent de dupliquer la mise en forme.

\section{Contrat de l'API}

L'API fournit un contrat HTTP entre le frontend et les traitements Python. Le
tableau \ref{tab:endpoints-api} résume les endpoints principaux.

\begin{longtable}{p{0.18\textwidth} p{0.28\textwidth} p{0.38\textwidth}}
\caption{Endpoints de l'API FastAPI}\label{tab:endpoints-api}\\
\toprule
Méthode & Route & Fonction \\
\midrule
\endfirsthead
\toprule
Méthode & Route & Fonction \\
\midrule
\endhead
GET & \texttt{/api/health} & Vérifier l'état du modèle et des données \\
GET & \texttt{/api/meta} & Retourner les bornes temporelles et la qualité \\
GET & \texttt{/api/overview} & Retourner les statistiques générales \\
GET & \texttt{/api/demand-analysis} & Retourner les agrégations d'analyse \\
GET & \texttt{/api/model-performance} & Retourner métriques et diagnostics \\
GET & \texttt{/api/forecast} & Comparer réel et prédit pour une date \\
POST & \texttt{/api/forecast/future} & Produire une prévision récursive \\
\bottomrule
\end{longtable}

Les requêtes sont validées avant le calcul. Une date inconnue, un horizon
inférieur à 1 ou supérieur à 168, ou des tableaux météo de longueur
incorrecte doivent produire une réponse d'erreur explicite. Cette validation
évite qu'une entrée invalide ne provoque une erreur difficile à interpréter
dans l'interface.

\section{Métriques d'évaluation}

La moyenne des erreurs absolues est définie par :

\begin{equation}
MAE = \frac{1}{n}\sum_{i=1}^{n}|y_i-\hat{y}_i|.
\end{equation}

La racine de l'erreur quadratique moyenne pénalise davantage les grandes
erreurs :

\begin{equation}
RMSE = \sqrt{\frac{1}{n}\sum_{i=1}^{n}(y_i-\hat{y}_i)^2}.
\end{equation}

Le coefficient $R^2$ compare la variance expliquée par le modèle à la variance
de la cible. Le WAPE est calculé par :

\begin{equation}
WAPE = 100 \times \frac{\sum_i |y_i-\hat{y}_i|}{\sum_i |y_i|}.
\end{equation}

Enfin, le biais est défini comme la moyenne de $\hat{y}_i-y_i$. Un biais
négatif correspond à une tendance à sous-estimer la demande. L'utilisation
conjointe de ces métriques évite de réduire l'analyse à un seul indicateur.

\section{Reproductibilité de l'expérimentation}

La reproductibilité dépend de quatre éléments : le fichier source, le code de
préparation, la configuration du modèle et la séparation temporelle. Le
notebook documente les étapes d'exploration et d'entraînement. Le backend
centralise les règles de features utilisées pendant l'inférence.

La graine aléatoire est fixée dans la configuration du modèle. Les features
sont ordonnées avant l'appel à XGBoost afin d'éviter qu'un changement de
l'ordre des colonnes ne produise une prédiction incohérente. Le modèle est
sauvegardé avec sa liste de variables et ses métriques dans un package Joblib.

\section{Description des sprints et itérations}

La réalisation a suivi les itérations suivantes :

\begin{enumerate}
	\item ingestion et contrôle de qualité du fichier de données ;
	\item exploration des tendances et création des variables temporelles ;
	\item entraînement et sélection du modèle XGBoost ;
	\item développement des endpoints FastAPI ;
	\item construction des pages React et des graphiques ;
	\item ajout des tests et vérification des prévisions récursives.
\end{enumerate}

\section{Implémentation des composants clés}

Le module d'ingestion convertit les timestamps, trie les observations et
impose une fréquence horaire. Les variables météorologiques sont propagées
vers l'avant lorsque cela est nécessaire, tandis que la demande cible n'est
pas imputée. Les lignes qui ne disposent pas des variables nécessaires sont
écartées de l'entraînement et de l'évaluation.

Le modèle utilise 32 variables : des variables calendaires, des variables
cycliques, des indicateurs de week-end et de jour férié, la température et
l'humidité, des retards de demande et des moyennes ou écarts-types glissants.
Les retards utilisés sont notamment de 1, 2, 3, 6, 12, 24, 48, 72, 168 et
336 heures.

\begin{figure}[H]
    \centering
    \includegraphics[width=1.1\textwidth]{perfocmance_2.jpeg}
    \caption{Performance du modèle}
    \label{fig:}
\end{figure}

\section{Résultats techniques}

Le modèle actif est un \texttt{XGBRegressor} utilisant l'objectif
\texttt{reg:squarederror}. Il comporte 1 328 arbres, un taux d'apprentissage de
0,02, une profondeur maximale de 5, un sous-échantillonnage des observations
de 0,85 et un sous-échantillonnage des variables de 0,85. La régularisation
L1 vaut 0,1 et la régularisation L2 vaut 3,0.

La sélection a porté sur huit configurations candidates. Elle s'est appuyée
sur des validations chronologiques en 2023, avec le WAPE moyen comme critère
principal et le RMSE comme critère secondaire. Le modèle retenu est
\texttt{enhanced_conservative}.

% Pour un projet d'IA : détailler l'entraînement et les métriques d'évaluation.
% Pour un projet de BI : détailler l'entrepôt de données et les tableaux de bord.

\section{Tests}

\subsection{Tests unitaires}

Les tests unitaires vérifient notamment la création des retards, des variables
cycliques et des statistiques glissantes. Ils contrôlent aussi les calculs de
métriques et la séparation chronologique des données.

\subsection{Tests d'intégration}

Les tests d'intégration couvrent l'état de santé de l'API, les métadonnées,
la vue d'ensemble, l'analyse de la demande, la performance du modèle et les
routes de prévision. Ils vérifient également les erreurs liées aux dates, aux
horizons et aux données météorologiques.

\subsection{Validation utilisateur}

La validation réalisée dans le dépôt est principalement technique et
automatisée. Elle vérifie que les pages frontend s'affichent, que les
métriques sont présentes et que l'appel de prévision future est correctement
effectué. Une validation métier par des utilisateurs finaux reste une étape
à réaliser avant toute utilisation opérationnelle.

\section{Guide d'utilisation de l'application}

L'application s'ouvre sur la vue d'ensemble. L'utilisateur y vérifie la
période couverte, le nombre d'observations et les principaux indicateurs. Il
peut ensuite accéder à l'analyse de la demande pour observer les profils
horaires et mensuels, comparer les jours ouvrés aux week-ends et examiner la
relation entre température et demande.

La page de performance présente la comparaison entre XGBoost et les
baselines. Elle permet de consulter les métriques globales, les erreurs, les
valeurs réelles et prédites ainsi que l'importance des variables. Cette page
sert de support à l'interprétation et non de preuve suffisante pour un
déploiement industriel.

La page Forecast propose deux usages. La consultation historique permet de
choisir une date existante et de comparer la valeur réelle à la valeur
prédite. La prévision future demande un horizon et, lorsque cela est
disponible, des tableaux météo. L'utilisateur doit lire l'avertissement sur
le caractère récursif des horizons longs.

\begin{enumerate}
\item Vérifier l'état de l'API dans la vue d'ensemble.
\item Contrôler les bornes temporelles et la qualité des données.
\item Explorer les profils de demande avant d'interpréter une prévision.
\item Comparer le modèle actif aux baselines.
\item Générer une prévision future avec un horizon raisonnable.
\item Conserver les captures et les paramètres utilisés pour la traçabilité.
\end{enumerate}

\section{Plan de validation fonctionnelle}

La validation fonctionnelle est organisée par parcours utilisateur. Chaque
parcours part d'une action et vérifie la réponse affichée. Le parcours santé
contrôle que l'API, les données et le modèle sont disponibles. Le parcours
analyse vérifie que les agrégations ne sont pas vides. Le parcours modèle
vérifie que les métriques et les baselines sont cohérentes. Le parcours
prévision vérifie le traitement d'une date historique puis d'un horizon futur.

\begin{longtable}{p{0.25\textwidth} p{0.35\textwidth} p{0.25\textwidth}}
\caption{Matrice de validation fonctionnelle}\label{tab:validation-fonctionnelle}\\
\toprule
Parcours & Action & Résultat attendu \\
\midrule
\endfirsthead
\toprule
Parcours & Action & Résultat attendu \\
\midrule
\endhead
Santé & Ouvrir la vue d'ensemble & Statut et bornes affichés \\
Analyse & Ouvrir Demand Analysis & Graphiques non vides \\
Performance & Ouvrir Model Performance & Métriques et baselines affichées \\
Historique & Choisir une date valide & Réel et prédit comparables \\
Future & Demander 1 à 168 heures & Série future retournée \\
Erreur & Envoyer un horizon invalide & Message d'erreur contrôlé \\
\bottomrule
\end{longtable}

\chapter{Résultats et Discussion}

\section{Présentation des résultats obtenus}

L'évaluation finale porte sur 8 666 observations communes de l'année 2024.
Le modèle actif est comparé à une baseline utilisant la valeur de la veille,
à une baseline utilisant la valeur de la même heure de la semaine précédente
et à un ancien modèle XGBoost.

\begin{table}[htbp]
	\centering
	\caption{Performances sur le test 2024}
	\label{tab:performances-2024}
	\begin{tabular}{lrrrrr}
		\toprule
		Modèle & MAE & RMSE & $R^2$ & WAPE & Biais \\
		\midrule
		Naïf 24 heures & 178,25 & 224,53 & 0,7860 & 6,09\% & -0,99 \\
		Naïf 168 heures & 144,33 & 181,99 & 0,8594 & 4,93\% & -1,15 \\
		Ancien XGBoost & 112,95 & 141,32 & 0,9152 & 3,86\% & -59,99 \\
		XGBoost actif & \textbf{103,72} & \textbf{130,22} & \textbf{0,9280} & \textbf{3,55\%} & \textbf{-28,58} \\
		\bottomrule
	\end{tabular}
	\\[0.2cm]
	\small Les erreurs MAE, RMSE et biais sont exprimées en MW.
\end{table}

\begin{figure}[H]
    \centering
    \includegraphics[width=1.1\textwidth]{prevision future_.jpeg}
    \caption{Wireframe Interface}
    \label{fig:}
\end{figure}

\section{Analyse et interprétation}

Le XGBoost actif obtient les meilleurs résultats parmi les modèles comparés.
Son WAPE de 3,55\% indique une erreur absolue pondérée inférieure à celles
des deux baselines et de l'ancien modèle. Le biais négatif de -28,58 MW montre
toutefois une tendance moyenne à la sous-estimation. Le modèle explique une
grande partie de la variabilité du test, avec un $R^2$ de 0,9280.

Ces résultats doivent être interprétés comme une évaluation historique à
court horizon avec des retards observés. Ils ne représentent pas directement
la précision d'une prévision future entièrement récursive.

\section{Comparaison avec l'existant et la baseline}

Par rapport à la baseline à 24 heures, le modèle actif réduit le MAE de
178,25 MW à 103,72 MW et le WAPE de 6,09\% à 3,55\%. Il améliore également la
baseline à 168 heures, dont le WAPE est de 4,93\%. L'ancien XGBoost atteint un
WAPE de 3,86\%, mais présente un biais plus important de -59,99 MW. Le modèle
actif offre donc le meilleur compromis observé sur le test 2024.

\section{Lecture détaillée des métriques}

La baseline à 24 heures fournit un premier repère avec un MAE de 178,25 MW.
Le modèle actif réduit cette erreur à 103,72 MW, soit une amélioration
absolue de 74,53 MW. Le RMSE passe de 224,53 MW à 130,22 MW, ce qui indique
une réduction importante des erreurs élevées. Le WAPE passe de 6,09\% à
3,55\%, ce qui facilite la comparaison entre périodes de demande différente.

La baseline hebdomadaire est plus compétitive : son MAE vaut 144,33 MW et son
WAPE 4,93\%. Cette observation confirme que la saisonnalité hebdomadaire est
importante. Malgré cela, le modèle enrichi conserve un avantage grâce à la
combinaison des retards, de la météo et des variables calendaires.

L'ancien XGBoost obtient un WAPE de 3,86\%, mais son biais de -59,99 MW est
plus défavorable que celui du modèle actif. Un modèle peut donc avoir une
bonne erreur moyenne tout en présentant une tendance systématique. L'analyse
des métriques doit tenir compte de cette différence.

\section{Analyse des erreurs par période}

Les erreurs sont analysées par heure, jour de la semaine, mois et niveau de
demande. Cette décomposition permet de repérer les périodes où le modèle est
moins stable. Les heures de pointe sont particulièrement importantes, car une
erreur à ces moments peut avoir un impact plus important sur la planification
que la même erreur pendant une période creuse.

La comparaison entre semaine et week-end permet de vérifier si les variables
de calendrier capturent suffisamment le changement de comportement. L'analyse
par mois permet de rechercher un effet saisonnier résiduel. Enfin, la
catégorisation en demande faible, normale et de pointe donne une lecture plus
opérationnelle des résultats.

\begin{figure}[H]
    \centering
    \includegraphics[width=1.0\textwidth]{analyse erreurs.jpeg}
    \caption{Analyse des erreurs par heure et par mois}
    \label{fig:}
\end{figure}

\section{Importance des variables}

L'importance des variables fournit un premier outil d'explicabilité. Les
retards de 24 et 168 heures sont attendus parmi les variables utiles, car ils
représentent les répétitions journalière et hebdomadaire. L'heure de la
journée décrit les changements de rythme. La température et les statistiques
glissantes ajoutent respectivement le contexte extérieur et le niveau récent
de la demande.

Cette analyse ne doit pas être interprétée comme une causalité physique
stricte. Une importance élevée signifie que la variable contribue aux
décisions des arbres dans le jeu étudié. Elle ne prouve pas que modifier cette
variable provoquerait directement une variation de la demande.

\begin{figure}[H]
    \centering
    \includegraphics[width=1.0\textwidth]{importance des variables.jpeg}
    \caption{Importance des variables}
    \label{fig:}
\end{figure}

\section{Intervalles et incertitude}

L'interface peut afficher des intervalles empiriques autour de certaines
prévisions. Ces intervalles sont estimés à partir du comportement historique
des erreurs ; ils ne constituent pas une garantie statistique. Pour obtenir
des intervalles calibrés, il faudrait étudier la distribution conditionnelle
des erreurs, réaliser une validation dédiée et mesurer la couverture sur une
période indépendante.

\section{Limites identifiées}

La première limite concerne les données : elles sont synthétiques et ne
permettent pas de conclure sur la précision d'un système électrique réel. Le
projet ne comporte pas non plus de fournisseur météo externe. En l'absence de
prévisions météorologiques, les valeurs utilisées pour l'avenir peuvent être
estimées à partir de moyennes historiques par mois et par heure.

La seconde limite concerne les prévisions multi-horizon. Elles sont
récursives, car les prédictions précédentes sont réinjectées dans les étapes
suivantes. Enfin, le dépôt ne fournit pas encore de déploiement de production,
de validation utilisateur formelle ou de garantie statistique pour les
intervalles empiriques affichés.

\chapter{Conclusion générale et Perspectives}

\section{Bilan du travail réalisé}

Ce projet a permis de réaliser un démonstrateur local complet de prévision
horaire de la demande électrique. La chaîne développée couvre l'ingestion des
données, le contrôle qualité, la création de variables, l'entraînement et
l'évaluation d'un modèle XGBoost, l'exposition des résultats par une API
FastAPI et leur consultation dans une interface React.

Les objectifs techniques ont été atteints : les données sont analysées selon
leur ordre chronologique, les performances sont comparées à des baselines et
une prévision récursive de 1 à 168 heures est disponible. Sur les données de
test 2024, le modèle actif obtient un WAPE de 3,55\% et un $R^2$ de 0,9280.
Ces résultats restent ceux d'un prototype évalué sur des données synthétiques
et ne constituent pas une validation industrielle.

\section{Compétences acquises}

La réalisation de ce projet a permis de développer des compétences en :

\begin{itemize}
	\item nettoyage, contrôle qualité et analyse exploratoire de données ;
	\item création de variables pour les séries temporelles ;
	\item entraînement, sélection et évaluation de modèles de régression ;
	\item conception d'API avec FastAPI ;
	\item développement d'interfaces React et de visualisations interactives ;
	\item écriture de tests backend et frontend ;
	\item documentation et communication de résultats techniques.
\end{itemize}

\section{Améliorations futures et pistes de recherche}

Plusieurs améliorations peuvent prolonger ce travail. La première consiste à
remplacer les données synthétiques par des mesures électriques réelles et à
réaliser une validation avec les utilisateurs de l'entreprise. Il serait aussi
pertinent d'intégrer un fournisseur de prévisions météorologiques afin de ne
plus dépendre uniquement de moyennes historiques.

Le système pourrait également intégrer des événements particuliers tels que
le Ramadan, après validation d'une source de calendrier fiable, ainsi que des
modèles supplémentaires pour comparer les performances. Enfin, un déploiement
industriel nécessiterait une authentification, une supervision, une gestion
des versions du modèle, une journalisation et une étude de la stabilité des
prévisions sur des horizons longs.

\renewcommand{\bibname}{Bibliographie et Webographie}
\addcontentsline{toc}{chapter}{Bibliographie et Webographie}

\begin{thebibliography}{9}
	\bibitem{xgboost}
	Chen, T. et Guestrin, C., \og XGBoost: A Scalable Tree Boosting System\fg{},
	\textit{Proceedings of the 22nd ACM SIGKDD International Conference on
	Knowledge Discovery and Data Mining}, 2016.
	\bibitem{xgboost-docs}
	XGBoost Developers, \textit{XGBoost Documentation},
	\url{https://xgboost.readthedocs.io/}, consulté le 13 septembre 2026.
	\bibitem{fastapi-docs}
	FastAPI Developers, \textit{FastAPI Documentation},
	\url{https://fastapi.tiangolo.com/}, consulté le 13 septembre 2026.
	\bibitem{react-docs}
	Meta Open Source, \textit{React Documentation},
	\url{https://react.dev/}, consulté le 13 septembre 2026.
	\bibitem{pandas-docs}
	The pandas development team, \textit{pandas Documentation},
	\url{https://pandas.pydata.org/docs/}, consulté le 13 septembre 2026.
	\bibitem{sklearn-docs}
	Scikit-learn Developers, \textit{Scikit-learn User Guide},
	\url{https://scikit-learn.org/stable/user_guide.html}, consulté le 13 septembre 2026.
\end{thebibliography}

\appendix

\chapter{Extraits de code}

Les extraits les plus importants du projet concernent la création des
variables temporelles et historiques, le chargement du modèle Joblib, la
construction de la réponse de l'API et la boucle de prévision récursive. Le
code source complet est organisé dans les modules
\texttt{backend/ingestion.py}, \texttt{backend/features.py},
\texttt{backend/modeling.py} et \texttt{backend/main.py}.

\begin{verbatim}
# Schéma simplifié de la prévision récursive
historique = charger_historique()
pour chaque heure future :
	variables = construire_features(historique)
	prediction = modele.predict(variables)
	historique.ajouter(prediction)
\end{verbatim}

\chapter{Captures d'écran supplémentaires}

Cette annexe est réservée aux captures complémentaires de l'application, aux
résultats détaillés et aux vues qui ne sont pas présentées dans le corps du
rapport.

\screenshotpage{Page Overview}{Insérer la capture complète de la page Overview.}
\screenshotpage{Indicateurs de la vue d'ensemble}{Insérer un zoom sur les cartes
de période, de volume, de moyenne, de minimum et de maximum.}
\screenshotpage{Série temporelle journalière}{Insérer le graphique de l'évolution
de la demande journalière.}
\screenshotpage{Analyse par heure}{Insérer les graphiques de demande moyenne et
de distribution par heure.}
\screenshotpage{Analyse par mois}{Insérer les graphiques de saisonnalité
mensuelle.}
\screenshotpage{Semaine et week-end}{Insérer la comparaison entre jours ouvrés
et week-end.}
\screenshotpage{Température et demande}{Insérer le nuage de points et la tendance
température-demande.}
\screenshotpage{Matrice de corrélation}{Insérer la heatmap des corrélations.}
\screenshotpage{Performance globale}{Insérer la page Model Performance complète.}
\screenshotpage{Comparaison des baselines}{Insérer le tableau ou graphique de
comparaison entre les modèles.}
\screenshotpage{Valeurs réelles et prédites}{Insérer la série des valeurs réelles
et prédites sur la période de test.}
\screenshotpage{Distribution des erreurs}{Insérer l'histogramme des erreurs et
les statistiques associées.}
\screenshotpage{Importance des variables}{Insérer le classement des features.}
\screenshotpage{Diagnostic par horizon}{Insérer le diagnostic des scores selon
l'horizon de prévision.}
\screenshotpage{Prévision historique}{Insérer un exemple de comparaison pour une
date historique.}
\screenshotpage{Prévision future}{Insérer la page Forecast avec une prévision
future de plusieurs heures.}
\screenshotpage{Réponse de l'API}{Insérer une capture de la réponse JSON d'un
endpoint important.}
\screenshotpage{Documentation technique}{Insérer une capture de la documentation
interactive FastAPI si elle est disponible.}
\screenshotpage{Tests backend}{Insérer la sortie des tests Pytest.}
\screenshotpage{Tests frontend}{Insérer la sortie des tests Vitest.}

\chapter{Annexe technique : configuration et exécution}

Cette annexe rassemble les éléments nécessaires pour reproduire le projet.
Les commandes exactes peuvent dépendre du système d'exploitation, mais le
principe reste identique : installer les dépendances Python, lancer l'API,
installer les dépendances frontend puis démarrer Vite.

\begin{verbatim}
python -m venv .venv
.venv\Scripts\activate
python -m pip install -r requirements.txt
uvicorn backend.main:app --reload

cd frontend
npm install
npm run dev
\end{verbatim}

Le fichier de données doit être placé dans le chemin attendu ou indiqué par
la variable d'environnement prévue par le backend. Avant de démarrer le
frontend, il faut vérifier que l'URL de l'API correspond à la configuration
locale. Le notebook peut ensuite être ouvert avec le même environnement Python
afin de reproduire l'exploration et l'entraînement.

\section{Checklist de reproduction}

\begin{enumerate}
\item Vérifier la version de Python et l'environnement virtuel.
\item Installer les dépendances inscrites dans \texttt{requirements.txt}.
\item Vérifier la présence du fichier de données et ses quatre colonnes.
\item Exécuter les contrôles d'ingestion.
\item Recalculer les features sans mélanger futur et passé.
\item Vérifier la séparation 2022, 2023 et 2024.
\item Entraîner ou charger l'artefact XGBoost.
\item Lancer les tests backend et frontend.
\item Démarrer l'API puis le tableau de bord React.
\item Conserver les versions, métriques et captures dans le dossier du rapport.
\end{enumerate}

\section{Structure logique du dépôt}

\begin{longtable}{p{0.3\textwidth} p{0.55\textwidth}}
\caption{Principaux éléments du dépôt}\label{tab:structure-depot}\\
\toprule
Élément & Rôle \\
\midrule
\endfirsthead
\toprule
Élément & Rôle \\
\midrule
\endhead
\texttt{backend/ingestion.py} & Chargement et contrôle des données \\
\texttt{backend/features.py} & Features causales et calendrier \\
\texttt{backend/modeling.py} & Entraînement et évaluation \\
\texttt{backend/main.py} & Routes FastAPI \\
\texttt{backend/tests/} & Tests du backend \\
\texttt{frontend/src/} & Application React et pages \\
\texttt{models/} & Artefacts et métadonnées du modèle \\
\texttt{notebook.ipynb} & Workflow d'exploration et d'expérimentation \\
\bottomrule
\end{longtable}

\chapter{Glossaire}

\begin{description}
	\item[API] Interface permettant à des applications de communiquer entre elles.
	\item[Baseline] Modèle de référence utilisé pour comparer une nouvelle solution.
	\item[Feature] Variable fournie au modèle d'apprentissage automatique.
	\item[MAE] Erreur absolue moyenne entre les valeurs prédites et réelles.
	\item[RMSE] Racine de l'erreur quadratique moyenne.
	\item[WAPE] Erreur absolue pondérée par la demande observée.
	\item[Prévision récursive] Prévision dans laquelle les prédictions précédentes
	servent de variables pour les étapes suivantes.
	\item[XGBoost] Bibliothèque de gradient boosting utilisée ici pour la régression.
\end{description}

\end{document}
